\documentclass[UTF8]{ctexart}
\usepackage{amsmath,amssymb}
\begin{document}
\section*{区域协同交通信号控制：完整当前问题定义}
本问题以实例给定的中观道路网络、车辆需求和信号相位为输入，优化区域协同交通信号控制策略。

\subsection*{环境与状态}
有向路网及智能体区域由输入给出。路段$e$具有储存上限$C_e$、车道数$n_e$、
长度$L_e$、自由行驶时间$\tau_e$和绿灯相位编号。相位预设为4种互斥流向。
完整环境状态$X_t$包括各路段FIFO车辆队列、在途车辆及剩余行驶时间、
车辆的固定最短距离路径及路径位置、未入网起点队列、当前相位、绿灯起始时刻、
黄灯/全红计时器和时钟。该完整状态具有马尔可夫转移，下述观测是完整状态的压缩表示。

起点$o$每秒产生泊松到达，率为
\[
\lambda_o(t)=\lambda_0m_o
\frac{1+0.23\sin(2\pi t/H+0.27o)+0.13\sin(4\pi t/H+0.6)}
{H^{-1}\sum_{u=0}^{H-1}[1+0.23\sin(2\pi u/H+0.27o)+0.13\sin(4\pi u/H+0.6)]}.
\]
目的地在其他节点中均匀抽取，路径按路段长度最短路确定。
未能入网的车辆在起点FIFO队列保留，不丢弃需求。
每秒依次产生需求、按首段剩余容量准入、将行驶结束车辆加入路段FIFO队列、
更新信号、按$(t+j)\bmod |V|$轮转次序处理路口放行。绿灯路段每两秒每车道
最多放行1辆；下游储存不足时阻塞，驶出车辆进入下一段并等待$\tau_e$秒，
或完成行程。必须始终满足
\[
0\le O_e(t)\le C_e,\quad O_e(t)=Q_e(t)+M_e(t),\qquad
G(t)=D(t)+\sum_eO_e(t)+B(t).
\]
$Q,M,O$分别为排队、在途和占用车辆数，$G,D,B$为累计产生、完成和起点积压数。
最后尚未完成的行程和起点积压全部报告。

\subsection*{策略与信号可行性}
每隔$\delta$秒每个区域选动作$a_i\in\{0,1\}$。动作0按
$\arg\max_{e\in\operatorname{in}(v)}O_e/(L_en_e)$选择预设绿灯相位，
对应原文Algorithm 1的密度及预设相位选择；动作1按周期顺序轮转相位。
切换仅在当前绿灯达到最小时长后执行；达到最大绿灯时长必须换至其他相位。
黄灯和全红持续输入给定秒数，期间不放行，结束后新绿灯计时从0开始。
仿真终点截断的最后一段绿灯可短于最小时长，不能超过最大时长。

智能体$i$的观测由总排队数三级分类（阈值5、10）、排队最大的两个路口编号、
两路口分类和上一动作共6项组成；与按区域中心距离选定的邻居观测拼接。
邻居平均动作$\bar a_i=|\mathcal N_i|^{-1}\sum_{j\in\mathcal N_i}a_j$用于策略；
实现中用邻居动作之和作为等价的表索引，因每个区域邻居数固定。
同一候选策略根据当前观测及以往动作选择信号控制。

\subsection*{目标与求解算法}
对实际评估的$E$条独立需求样本轨迹，目标为
\[
\min_\pi J_E(\pi)=\frac1{EAT}\sum_{e=1}^E\sum_{t=1}^T
\sum_{i=1}^A\sum_{v\in G_i}\sum_{l\in\operatorname{in}(v)}Q_l^{e,t}.
\]
这是已记录轨迹的实际平均路内排队数。起点未入网等待不在这个目标内，
另行报告。
求解采用表格Q学习作为寻找可行策略的启发式，奖励为
\[
f_i^t=\frac1{|G_i|}\sum_{v\in G_i}
\frac{\sum_{l\in\operatorname{in}(v)}M_l^t v_l}{\sum_{l\in\operatorname{in}(v)}L_ln_l},
\quad r_i^t=\operatorname{sgn}(f_i^t-f_i^{t-1})
+\sum_{j\in\mathcal N_i}\frac{\operatorname{sgn}(f_j^t-f_j^{t-1})}{d_{ij}}.
\]
\[
Q_i\leftarrow(1-\alpha)Q_i+\alpha
\left[r_i+\gamma\max_{a'}Q_i(X',a',\bar a_i')\right],\qquad
\pi_i(a)\propto\exp(+\beta Q_i(a)).
\]
终止步去掉未来项，正号对应奖励最大化方向。实现将二动作logit裁剪到$[-60,60]$
以避免浮点溢出，实际概率按裁剪值定义。每个场景预算3600秒，按控制步检查时限。
训练次数与评估次数是输入算法参数。
输出为可行控制策略及其轨迹目标。
\end{document}
